% Part: incompleteness
% Chapter: incompleteness-provability
% Section: second-incompleteness-thm

\documentclass[../../../include/open-logic-section]{subfiles}

\begin{document}

\olfileid[hi]{inc}{inp}{2in}

\olsection{द्वितीय अपूर्णता-प्रमेय}

इस कथन को हम कैसे व्यक्त कर सकते हैं कि $\Th{PA}$ अपनी संगतता स्वयं सिद्ध
नहीं करता? $\Th{PA}$ के असंगत होने का अर्थ यह कहना है कि
$\Th{PA} \Proves \eq[0][1]$। अतः हम संगतता-कथन $\OCon[\Th{PA}]$ को
!!{sentence} $\lnot \OProv[\Th{PA}](\gn{\eq[0][1]})$ मान सकते हैं; तब
अगला प्रमेय वांछित परिणाम देता है:

\begin{thm}
\ollabel{thm:second-incompleteness}
यदि $\Th{PA}$ संगत है, तो $\Th{PA}$,~$\OCon[\Th{PA}]$ को !!{derive} नहीं
करता।
\end{thm}

यह ध्यान देना महत्त्वपूर्ण है कि प्रमेय $\OCon[\Th{PA}]$ के विशिष्ट
निरूपण (अर्थात् $\OProv[\Th{PA}](y)$ के विशिष्ट निरूपण) पर निर्भर है। हम
केवल यह उपयोग करेंगे कि $\OProv[\Th{PA}](y)$ का निरूपण तीनों
!!{derivability} प्रतिबंधों को संतुष्ट करता है; अतः प्रमेय ऐसे
!!a{derivability} विधेय वाले प्रत्येक सिद्धांत तक व्यापक होता है जिसमें ये
गुण हों।

ग्योडेल (G\"odel) के तर्क की रूपरेखा पढ़ना ज्ञानवर्धक है, क्योंकि प्रमेय
एक प्रभावी अंतिम पंक्ति की तरह निकलता है। तर्क ऐसा है। $!G_\Th{PA}$ को वह
ग्योडेल (G\"odel) वाक्य मानें जिसे हमने \olref[1in]{thm:first-incompleteness} के
प्रमाण में बनाया था। हमने दिखाया है: ``यदि $\Th{PA}$ संगत है, तो
$\Th{PA}$,~$!G_\Th{PA}$ को !!{derive} नहीं करता।'' यदि इसे $\Th{PA}$
\emph{के भीतर} औपचारिक बनाएँ, तो हमें इसका प्रमाण मिलता है:
\[
\OCon[\Th{PA}] \lif \lnot \OProv[\Th{PA}](\gn{!G_\Th{PA}}).
\]
अब मान लें $\Th{PA}$,~$\OCon[\Th{PA}]$ को !!{derive} करता है। तब वह
$\lnot \Prov[\Th{PA}](\gn{!G_\Th{PA}})$ को !!{derive} करता है। किंतु
चूँकि $!G_\Th{PA}$ ग्योडेल (G\"odel) वाक्य है, यह $!G_\Th{PA}$ के समतुल्य है। अतः
$\Th{PA}$,~$!G_\Th{PA}$ को !!{derive} करता है।

किंतु हम जानते हैं कि यदि $\Th{PA}$ संगत है, तो वह $!G_\Th{PA}$ को
!!{derive} नहीं करता!{} अतः यदि $\Th{PA}$ संगत है, तो वह
$\OCon[\Th{PA}]$ को !!{derive} नहीं कर सकता।

तर्क को अधिक सटीक बनाने के लिए हम $!G_\Th{PA}$ को~$\Th{PA}$ का ग्योडेल (G\"odel)
वाक्य लेंगे और !!{derivability} प्रतिबंधों (P1)--(P3) का उपयोग करके दिखाएँगे
कि $\Th{PA}$,~$\OCon[\Th{PA}] \lif !G_\Th{PA}$ को !!{derive} करता है।
इससे दिखेगा कि $\Th{PA}$,~$\OCon[\Th{PA}]$ को !!{derive} नहीं करता। यहाँ
$\Th{PA}$ के भीतर प्रमाण की रूपरेखा है। (सरलता के लिए हम $\Th{PA}$ के
अधोलेख हटा देते हैं।)
\begin{align}
& !G \liff \lnot \OProv(\gn{!G}) \ollabel{G2-1}\\
& \qquad\text{$!G$ ग्योडेल (G\"odel) वाक्य है}\notag \\
& !G \lif \lnot \OProv(\gn{!G}) \ollabel{G2-2}\\
  & \qquad\text{\olref{G2-1} से} \notag\\
& !G \lif
  (\OProv(\gn{!G}) \lif \lfalse) \ollabel{G2-3}\\
  & \qquad\text{\olref{G2-2} से, तर्क द्वारा}\notag\\
& \OProv(\gn{
    !G \lif
    (\OProv(\gn{!G}) \lif \lfalse)
  }) \ollabel{G2-4}\\
  & \qquad\text{\olref{G2-3} से, प्रतिबंध P1 द्वारा} \notag\\
& \OProv(\gn{!G}) \lif
  \OProv(\gn{
    (\OProv(\gn{!G}) \lif \lfalse)
    }) \ollabel{G2-5}\\
  & \qquad\text{\olref{G2-4} से, प्रतिबंध P2 द्वारा} \notag\\
& \OProv(\gn{!G}) \lif (\OProv(\gn{\OProv(\gn{!G})}) \lif \OProv(\gn{\lfalse})) \ollabel{G2-6}\\
  & \qquad\text{\olref{G2-5} से, प्रतिबंध P2 और तर्क द्वारा} \notag\\
& \OProv(\gn{!G}) \lif
  \OProv(\gn{\OProv(\gn{!G})}) \ollabel{G2-7}\\
   & \qquad\text{P3 द्वारा} \notag\\
& \OProv(\gn{!G}) \lif \OProv(\gn{\lfalse}) \ollabel{G2-8}\\
  & \qquad \text{\olref{G2-6} और \olref{G2-7} से, तर्क द्वारा}\notag\\
& \OCon \lif \lnot \OProv(\gn{!G}) \ollabel{G2-9}\\
  & \qquad\text{\olref{G2-8} का प्रतिधनात्मक रूप और
  $\OCon \ident \lnot \OProv(\gn{\lfalse})$}\notag \\
& \OCon \lif !G \notag\\
  & \qquad\text{\olref{G2-1} और \olref{G2-9} से, तर्क द्वारा।}\notag
\end{align}
ऊपर तर्क का उपयोग प्रतिज्ञप्तिक तर्कशास्त्र के केवल प्राथमिक तथ्यों के लिए
है; उदाहरणतः \olref{G2-3},
$\Proves \lnot!A \liff (!A\lif \lfalse)$ का उपयोग करता है और
\olref{G2-8},
$!A \lif (!B \lif !C), !A \lif !B \Proves !A \lif !C$ का।
\olref{G2-5} और \olref{G2-6} में प्रतिबंध~P2 का उपयोग P2 के उदाहरणों
$\OProv(\gn{!A \lif !B}) \lif
(\OProv(\gn{!A}) \lif \OProv(\gn{!B}))$ पर निर्भर करता है। पहले में
$!A \ident !G$ और $!B \ident \OProv(\gn{!G}) \lif \lfalse$; दूसरे में
$!A \ident \OProv(\gn{G})$ और $!B \ident \lfalse$।

द्वितीय अपूर्णता-प्रमेय का अधिक अमूर्त रूप यह है:

\begin{thm}
\ollabel{thm:second-incompleteness-gen} $\Th{T}$,~$\Th{Q}$ का कोई संगत,
!!{axiomatized} विस्तार हो और $\OProv[\Th{T}](y)$ ऐसा कोई सूत्र हो जो
$\Th{T}$ के लिए !!{derivability} प्रतिबंध P1--P3 को संतुष्ट करता है। तब
$\Th{T}$,~$\OCon[T]$ को !!{derive} नहीं करता।
\end{thm}

\begin{prob}
दिखाएँ कि $\Th{PA}$,~$!G_{\Th{PA}} \lif \OCon[\Th{PA}]$ को !!{derive}
करता है।
\end{prob}

\begin{digress}
कथा का निष्कर्ष यह है कि गणित का कोई ``उचित'' संगत सिद्धांत अपने ही
संगतता-कथन को !!{derive} नहीं कर सकता। मान लें $\Th{T}$ गणित का ऐसा
सिद्धांत है जिसमें $\Th{Q}$ और हिल्बर्ट का ``परिमितीय'' तर्क (वह चाहे जो
भी हो) सम्मिलित हैं। तब पूरा $\Th{T}$,~$\Th{T}$ के संगतता-कथन को
!!{derive} नहीं कर सकता; अतः और भी स्पष्टतः, उसका परिमितीय खंड भी
$\Th{T}$ के संगतता-कथन को !!{derive} नहीं कर सकता। इस अर्थ में ``समूचे
गणित'' का कोई परिमितीय संगतता-प्रमाण नहीं हो सकता।

``परिमितीय'' पद की व्याख्या में कुछ छूट है, और 1931 के शोधपत्र में ग्योडेल (G\"odel)
इस संभावना को स्वीकार करते हैं कि जिसे हम ``परिमितीय'' मानें वह उस प्रकार
के गणित से बाहर हो सकता है जिसे हिल्बर्ट औपचारिक बनाना चाहते थे। किंतु
ग्योडेल (G\"odel) उदार हो रहे थे; आज यह देखना कठिन है कि हम ऐसी कोई चीज़ कैसे पा
सकते हैं जिसे उचित रूप से परिमितीय कहा जा सके, किंतु जिसे, मान लें,
चयन-अभिगृहीत सहित ज़र्मेलो--फ्रेंकेल समुच्चय-सिद्धांत $\Th{ZFC}$ में
औपचारिक न बनाया जा सके।
\end{digress}

\end{document}
